% Préambule commun aux manuels de Tutti (français et anglais).
% Compilation : pdflatex manuel-fr.tex (deux fois), idem manuel-en.tex, depuis docs/src/.
\documentclass[11pt,a4paper]{article}
\usepackage[utf8]{inputenc}
\usepackage[T1]{fontenc}
\usepackage{lmodern}
\usepackage{microtype}
\usepackage[a4paper,margin=2cm,headheight=24pt]{geometry}
\usepackage{graphicx}
\usepackage[table]{xcolor}
\usepackage[most]{tcolorbox}
\usepackage{enumitem}
\usepackage{booktabs}
\usepackage{array}
\usepackage{tabularx}
\usepackage[nobottomtitles*]{titlesec}
\usepackage{needspace}
\usepackage{fancyhdr}
\usepackage{fontawesome5}
\usepackage{etoolbox}
\DeclareUnicodeCharacter{2192}{\ensuremath{\rightarrow}}
\DeclareUnicodeCharacter{2191}{\ensuremath{\uparrow}}
\DeclareUnicodeCharacter{2193}{\ensuremath{\downarrow}}
\DeclareUnicodeCharacter{25B2}{\ensuremath{\blacktriangle}}
\DeclareUnicodeCharacter{25BC}{\ensuremath{\blacktriangledown}}
\DeclareUnicodeCharacter{21BA}{\ensuremath{\circlearrowleft}}
\DeclareUnicodeCharacter{21BB}{\ensuremath{\circlearrowright}}
\usepackage{amssymb}

% Palette de l'application : ivoire, bordeaux, or
\definecolor{bordeaux}{HTML}{2C6E49}
\definecolor{bordeauxl}{HTML}{E4F0E8}
\definecolor{or}{HTML}{B8860B}
\definecolor{orl}{HTML}{FBF3DC}
\definecolor{sapin}{HTML}{2C7A55}
\definecolor{sapinl}{HTML}{E2F1E9}
\definecolor{sky}{HTML}{1C6DB4}
\definecolor{skyl}{HTML}{E6F0F8}
\definecolor{summit}{HTML}{6B7A8F}
\definecolor{night}{HTML}{2A2420}
\definecolor{ink}{HTML}{1F2430}
\AtBeginDocument{\color{ink}}

\titleformat{\section}
  {\Large\bfseries\sffamily\color{bordeaux}}
  {\makebox[1.6em][l]{\thesection}}{0.3em}{}
  [{\vspace{1pt}\color{bordeaux!35}\titlerule[1.3pt]}]
\titlespacing*{\section}{0pt}{14pt}{7pt}
\titleformat{\subsection}{\large\bfseries\sffamily\color{night}}{\thesubsection}{0.5em}{}
\pretocmd{\section}{\needspace{10\baselineskip}}{}{}
\pretocmd{\subsection}{\needspace{7\baselineskip}}{}{}
\widowpenalty=10000
\clubpenalty=10000
\raggedbottom
\setlength{\parindent}{0pt}
\setlength{\parskip}{6pt}
\setlist[enumerate]{itemsep=3pt,leftmargin=6mm}
\setlist[itemize]{itemsep=2pt,leftmargin=5mm}

\pagestyle{fancy}
\fancyhf{}
\renewcommand{\headrulewidth}{0.4pt}
\renewcommand{\headrule}{\hbox to\headwidth{\color{bordeaux!40}\leaders\hrule height \headrulewidth\hfill}}
\fancyfoot[C]{\sffamily\small\color{summit}\thepage}
\renewcommand{\sectionmark}[1]{\markboth{#1}{}}

\newcommand{\menu}[1]{\textsf{\textbf{\textcolor{bordeaux}{#1}}}}
\newcommand{\touche}[1]{\tcbox[on line,boxsep=0pt,left=2pt,right=2pt,top=1.5pt,bottom=1.5pt,
  colback=white,colframe=summit,boxrule=0.5pt,arc=1.5pt]{\footnotesize\textsf{#1}}}
\newcommand{\fichier}[1]{\texttt{\small #1}}

\newcommand{\boite}[5]{\newtcolorbox{#1}{enhanced,unbreakable,colback=#2,colframe=#3,
  coltitle=white,fonttitle=\bfseries\sffamily\small,boxrule=0.9pt,arc=2.5pt,
  left=2.5mm,right=2.5mm,top=1.2mm,bottom=1.2mm,title={#4~~#5}}}
\newcommand{\capture}[2][1]{\par\needspace{8\baselineskip}\begin{center}\setlength{\fboxsep}{0pt}\setlength{\fboxrule}{0.6pt}%
  \fcolorbox{summit!50}{white}{\includegraphics[width=\dimexpr#1\linewidth-1.3pt\relax]{#2}}\end{center}}

\usepackage[english]{babel}
\usepackage[colorlinks=true,linkcolor=bordeaux,urlcolor=bordeaux]{hyperref}
\hypersetup{pdftitle={Da Capo — user guide},pdfauthor={},pdfcreator={Da Capo}}
\boite{astucebox}{orl}{or!85!black}{\faLightbulb}{Good to know}
\boite{securitebox}{sapinl}{sapin}{\faShield*}{Your scores are safe}
\boite{attentionbox}{alertl}{alert}{\faExclamationTriangle}{Caution}
\newcommand{\astuce}[1]{\par\needspace{5\baselineskip}\begin{astucebox}#1\end{astucebox}}
\newcommand{\attention}[1]{\par\needspace{5\baselineskip}\begin{attentionbox}#1\end{attentionbox}}
\definecolor{alert}{HTML}{B5382B}
\definecolor{alertl}{HTML}{F8E6E3}
\fancyhead[L]{\sffamily\small\color{bordeaux}\textbf{DA CAPO — USER GUIDE}}
\fancyhead[R]{\sffamily\small\color{summit}\leftmark}
\input{version}

\begin{document}

\begin{titlepage}
\centering
\vspace*{2.2cm}
\includegraphics[width=3.6cm]{../../logo_1024.png}\par
\vspace{8mm}
{\sffamily\Huge\bfseries\color{bordeaux} DA CAPO\par}
\vspace{3mm}
{\sffamily\LARGE\color{ink} Get your scores out of Finale\par}
\vspace{2mm}
{\sffamily\large\color{summit} User guide \textbullet{} version \versiondacapo\par}
\vspace{12mm}
\begin{tcolorbox}[enhanced,width=0.82\linewidth,colback=orl,colframe=or,arc=4pt,boxrule=1.2pt,left=5mm,right=5mm,top=3mm,bottom=3mm]
\sffamily\centering
\textit{Da capo}: “from the beginning”, in music.\\[2mm]
\raggedright
\textbf{1.} Score folders \quad \textbf{2.} Working folder \quad \textbf{3.} Finale \quad \textbf{4.} Start
\end{tcolorbox}
\vfill
{\footnotesize\color{summit}\sffamily Windows \textbullet{} macOS \textbullet{} Linux\par}
\end{titlepage}

\tableofcontents
\newpage

\section{Why Da Capo}

Finale has not been developed since 2024. It still works, but no update will ever repair it if
Windows breaks it one day. Its \fichier{.mus} and \fichier{.musx} files \textbf{only open in
Finale}. The way out is called \textbf{MusicXML}: an open format read by MuseScore Studio (free),
Dorico and Sibelius.

Finale can export a score to MusicXML, but one at a time. With hundreds of scores, that means days of
clicking. \textbf{Da Capo does the clicking for you}, score after score, while you do something
else — and checks the result.

\begin{center}
\renewcommand{\arraystretch}{1.35}
\begin{tabularx}{\linewidth}{>{\bfseries\color{bordeaux}}l X}
\toprule
1. Find & every Finale score in the chosen folders, subfolders included. \\
2. Copy & each one into the working folder, copy checked byte for byte. \\
3. Export & Finale opens each \emph{copy}, exports it to MusicXML (and a reference PDF), closes it. \\
4. Check & every MusicXML is read back: parts, bars, notes, lyrics, dynamics. \\
5. Prepare & every MusicXML becomes a MuseScore score (\fichier{.mscz}) + a PDF to compare. \\
6. Report & \fichier{rapport.html}: what worked, what failed, what deserves a look. \\
\bottomrule
\end{tabularx}
\end{center}

\begin{securitebox}
\begin{itemize}
\item The original files are \textbf{never opened or modified}: Finale only sees copies.
\item The copies are checked (SHA-256 fingerprint) and made read-only.
\item Nothing is sent over the Internet: everything happens on your computer.
\end{itemize}
\end{securitebox}

\section{Before you start}

\begin{enumerate}
\item \textbf{Install Da Capo}: \fichier{DaCapo-Setup.exe} on Windows (no administrator password).
      On first launch, Windows may show “Windows protected your PC”: \menu{More info} →
      \menu{Run anyway}.
\item \textbf{Install MuseScore Studio} (optional but recommended) from \fichier{musescore.org}
      — not \fichier{musescore.com}, a partly paid sheet-music site. The software is free.
\item \textbf{Open the Finale that Da Capo will use once by hand}, especially if it is never used:
      answer what it asks on first launch (authorisation, update, welcome window), then
      \textbf{close it}.
\item \textbf{Before Finale 25}, the reference PDF goes through the free virtual printer
      \textbf{CutePDF Writer}: check that it is installed.
\end{enumerate}

\astuce{If several Finale versions are installed, Da Capo suggests the most recent: it opens files
from older versions without trouble and writes a much more faithful MusicXML (version 4.0 for Finale
27, versus a “basic MusicXML” for Finale 2003 to 2010).}

\section{The four steps}

\capture[0.78]{../capture_prete-en.png}

\begin{description}[style=nextline,leftmargin=5mm,font=\sffamily\bfseries\color{bordeaux}]
\item[1. Where are your Finale scores?]
\menu{Add a folder…}, or drag the folder from Explorer. All subfolders are searched. Da Capo
immediately counts the scores found (\fichier{.mus} up to Finale 2011, \fichier{.musx} after).
Finale's backup copies (“Backup of…”) are ignored.
\item[2. Where to put the result?]
An empty folder, on a disk with room to spare (allow three to four times the size of the scores).
Everything goes there, in the same tree as your original folders.
\item[3. With which Finale?]
The most recent is suggested and recommended. The first time, click \menu{Diagnostics…}
(section~\ref{sec:diagnostic}).
\item[4. What to produce?]
\textbf{MusicXML}, always. The \textbf{reference PDF}: the score exactly as Finale prints it — your
yardstick for checking, and an archive that will always open. \textbf{MuseScore files}, if MuseScore
Studio is installed.
\end{description}

Then click \menu{Start}.

\section{During the export}

\capture[0.78]{../capture_export-en.png}

Finale opens and works on its own: you can see it open each score, export, close. \textbf{Do not
touch the mouse or keyboard}: a click at the wrong moment would divert a dialog box. Allow from a few
seconds to a minute per score; for several hundred, start in the evening and let it run overnight.
The computer does not go to sleep.

\begin{itemize}
\item \menu{Stop} finishes the current score, then stops cleanly.
\item The button then becomes \menu{Resume}: Da Capo skips everything already done. Even after a
      power cut, it only redoes the score that was in progress.
\item If Finale shows a message (old file to convert, missing font…), Da Capo reads it, records it in
      the report, then closes it. To “Save changes?”, it always answers \textbf{No}.
\item If Finale stops responding, Da Capo stops it, restarts it and moves on to the next score; the
      stubborn score is marked as failed and will be retried on resume.
\end{itemize}

\section{The working folder}

\begin{center}
\renewcommand{\arraystretch}{1.3}
\begin{tabularx}{\linewidth}{>{\ttfamily\small}l X}
\toprule
originaux/ & the copies of your scores, read-only \\
musicxml/ & the MusicXML files exported by Finale — keep them safe \\
pdf-finale/ & the reference PDFs \\
musescore/ & the MuseScore scores (\fichier{.mscz}) and their PDF \\
rapport.html & the report, to open in a web browser \\
\_dacapo/ & the log used to resume, and the diagnostics record \\
\bottomrule
\end{tabularx}
\end{center}

The whole folder can be moved or copied (to a backup disk, for example): the report's links follow.

\section{The report}

\menu{Open the report} shows a page with, at the top, the summary (scores, MusicXML, PDF, MuseScore,
failures), then one line per score:

\begin{center}
\renewcommand{\arraystretch}{1.35}
\begin{tabularx}{\linewidth}{l X}
\toprule
\textcolor{alert}{\rule{4pt}{9pt}}\ \textbf{red} & the export failed; the reason is written (and Finale's messages) \\
\textcolor{or}{\rule{4pt}{9pt}}\ \textbf{orange} & exported, but worth checking: empty score, Finale message during the export \\
\textcolor{sapin}{\rule{4pt}{9pt}}\ \textbf{green} & all went well \\
\bottomrule
\end{tabularx}
\end{center}

For each score: the number of parts, bars and notes found in the MusicXML, and links to every file
produced.

\section{Checking a score}

Open the \textbf{reference PDF} (Finale) and the \textbf{MuseScore PDF} side by side. Notes,
rhythms, lyrics, dynamics and parts come through MusicXML well. What often needs touching up in
MuseScore:
\begin{itemize}
\item the \textbf{layout}: system and page breaks, spacing, staff size;
\item \textbf{custom expressions} and rare symbols, special music fonts;
\item the \textbf{shape of slurs} and the fine position of text;
\item anything that only concerns Finale's audio playback (Human Playback).
\end{itemize}
Start with the scores you still use; the others stay archived in three copies (original, MusicXML,
PDF).

\attention{Keep the original \fichier{.mus} files and the copies in the working folder. One day,
a better converter may redo the export from them.}

\section{Diagnostics}\label{sec:diagnostic}

Each Finale version arranges its commands a little differently. Da Capo finds them by name, in
English and in French. The \menu{Diagnostics…} button checks this on your computer: it opens Finale
with a copy of one score, records its menus and windows, then closes everything, without changing
anything.

\capture[0.78]{../capture_diagnostic-en.png}

If everything is recognised, you are ready. Otherwise, \menu{Send the record} sends it to the
developer, who will adapt Da Capo to your Finale; the record is also saved in \fichier{\_dacapo}.

\section{Common problems}

\begin{center}
\renewcommand{\arraystretch}{1.35}
\begin{tabularx}{\linewidth}{>{\raggedright\arraybackslash}p{0.38\linewidth} X}
\toprule
\textbf{Message} & \textbf{What to do} \\
\midrule
Finale does not start & Open it once by hand (authorisation, update), close it, try again. \\
Command not found in Finale's menus & Run a diagnostic and send it. \\
CutePDF Writer printer not found & Install it (free), or untick the reference PDF. \\
Finale did not open the score in time & Damaged or very large file: open it by hand in Finale to see. \\
MuseScore took too long & Open the MusicXML directly in MuseScore. \\
\bottomrule
\end{tabularx}
\end{center}

\section{Updates and privacy}

At startup, Da Capo quietly checks for a new version and offers to install it. With your consent, a
report is sent if it runs into a problem: program and system versions, the error; your file names
and user name are hidden. No score is ever sent.

On macOS and Linux, Finale cannot be driven: Da Capo copies the scores, checks MusicXML files
exported elsewhere and prepares the MuseScore files.

\vfill
{\footnotesize\color{summit}\sffamily
Finale is a trademark of MakeMusic, MuseScore of Muse Group; Da Capo is affiliated with neither.
Fonts Source Sans 3 and EB Garamond (SIL Open Font License). Examples in this guide use fictitious files.\par}

\end{document}
